% ============================================================
% SEÇÃO 7 – MODELO ARQUITETURAL
% ============================================================

\section{MODELO ARQUITETURAL}

A arquitetura do LifeCity é estruturada em três camadas bem definidas,
seguindo o padrão de separação de responsabilidades:

\begin{figure}[H]
  \centering
  \includegraphics[width=0.6\textwidth]{images/arquitetura.png}
  \caption{Modelo arquitetural do LifeCity}
\end{figure}

\begin{itemize}
  \item \textbf{Camada de apresentação — Flutter App:} Responsável pela
        interface com o usuário. Desenvolvida em Flutter/Dart, comunica-se
        exclusivamente com a camada de aplicação via requisições HTTP REST,
        utilizando a biblioteca \texttt{Dio}. Gerencia autenticação local por
        meio de tokens JWT armazenados de forma segura no dispositivo.

  \item \textbf{Camada de aplicação — Node.js/Express (Backend):} Responsável
        por toda a lógica de negócio, autenticação e autorização. Implementa
        API RESTful com os seguintes grupos de endpoints:
        \begin{itemize}
          \item \texttt{/api/auth} — cadastro, login, logout, refresh token;
          \item \texttt{/api/users} — perfil, edição, XP, amizades;
          \item \texttt{/api/complaints} — CRUD de reclamações, curtidas, comentários, testemunhos;
          \item \texttt{/api/events} — CRUD de eventos;
          \item \texttt{/api/achievements} — conquistas do usuário;
          \item \texttt{/api/notifications} — listagem e marcação como lida;
          \item \texttt{/api/missions} — missões ativas, detalhes;
          \item \texttt{/api/missions/teams} — criar equipe, convidar, aceitar/recusar convite, editar equipe;
          \item \texttt{/api/missions/teams/:id/messages} — listar mensagens do chat da equipe (GET) e enviar nova mensagem (POST); acesso restrito a membros ativos;
          \item \texttt{/api/reports} — denúncias de reclamações e usuários com moderação automática por threshold;
          \item \texttt{/api/auth/forgot-password} e \texttt{/api/auth/reset-password} — recuperação de senha por e-mail com código de verificação de uso único.
        \end{itemize}
        Gerencia tokens JWT (access token de curta duração + refresh token de
        longa duração), incluindo renovação automática de sessão e invalidação
        no logout. Checadores de conquistas (\texttt{achievementChecker.js}),
        de progresso de missões (\texttt{missionProgressChecker.js}) e de
        confiabilidade do usuário (\texttt{cityValidator.js}) são disparados
        de forma assíncrona após ações do usuário. O envio de e-mails
        transacionais (recuperação de senha) é feito via \texttt{Nodemailer}
        integrado ao Gmail com autenticação por senha de aplicativo. O
        módulo \texttt{cityValidator.js} também é invocado de forma síncrona
        no cadastro para validar o CEP informado pelo usuário.

  \item \textbf{Camada de dados — PostgreSQL via Supabase:} Responsável pelo
        armazenamento persistente dos dados. O banco PostgreSQL é gerenciado
        pela plataforma Supabase, que também provê o serviço de armazenamento
        de arquivos de mídia (fotos de perfil e imagens de reclamações).
\end{itemize}

O aplicativo Flutter nunca acessa o banco de dados diretamente; toda
comunicação passa obrigatoriamente pelo backend Node.js, garantindo
controle de acesso e segurança centralizada.

\subsection*{Filtro inteligente de reclamações no mapa}

Para evitar a poluição visual do mapa em cidades com grande volume de
reclamações, o aplicativo aplica um filtro de relevância baseado na
distância entre cada reclamação e a posição GPS atual do usuário.
Reclamações a até 5\,km são exibidas normalmente (tamanho completo).
Reclamações além desse raio são exibidas apenas se possuírem
pontuação de relevância igual ou superior a 3, onde a pontuação é
calculada como a soma de curtidas, comentários e testemunhos
(\texttt{likesCount + commentsCount + witnessCount}).
Pins de alta relevância exibidos fora do raio próximo aparecem
levemente menores e com opacidade reduzida, comunicando visualmente
a diferença. O filtro é recalculado automaticamente sempre que a
posição do usuário é atualizada pelo GPS, sem necessidade de
interação. A distância é calculada pelo algoritmo de Haversine
diretamente no cliente Flutter, sem requisições adicionais ao backend.
